\section{Dimension-Free Hellinger Geometry}
\label{sec:rotational-fourier-defect-and-hellinger-coercivity}

The rotational Fourier defect distinguishes every non-Gaussian coordinate law
without dividing by its characteristic function.  Fixed Fourier frequencies
then produce bounded tests for all matrix directions, and the matrix Fisher
operator identifies their infinitesimal cost.  The goal of this section is the
dimension-free local estimate in
\cref{prop:local-fisher-hellinger}; the global estimate is derived from the
classification in \cref{sec:globalization-proof}.

\subsection{A product of centered complex exponentials}
\label{sec:product-of-centered-exponentials}

This subsection requires only $\E X=0$ and $\E|X|<\infty$.
Neither a density nor a second moment is needed.
Write
$\phi(u)=\E e^{iuX}$, and for independent copies $X_1,X_2$ put
\[
 O_\theta=\begin{pmatrix}\cos\theta&-\sin\theta\\
                         \sin\theta&\cos\theta\end{pmatrix},
 \qquad
 W_{\xi,\eta}(x,y)
   =(e^{i\xi x}-\phi(\xi))(e^{i\eta y}-\phi(\eta)).
\]
This fixes the sign convention for rotation.
Centering removes the two marginal responses.
It leaves only the dependence created by rotation.
In particular,
\begin{equation}
 \E W_{\xi,\eta}(X_1,X_2)=0,
 \qquad
 \E|W_{\xi,\eta}(X_1,X_2)|^2
 =(1-|\phi(\xi)|^2)(1-|\phi(\eta)|^2)\le1.
 \label{eq:centered-exponential-product-variance}
\end{equation}
Pointwise, $|W_{\xi,\eta}|\le4$, so its variance under an arbitrary
comparison law is at most $16$.  Both bounds are dimension-free and require
no tail assumption.

\subsection{Gaussian characterization by a global linear ODE}
\label{sec:fourier-ode}

The next lemma makes the rotational Fourier defect a logarithm-free rigidity
tool, valid even when \(\phi\) has zeros.  It is an infinitesimal version of
the Gaussian rigidity behind independence characterizations
\cite{kac1939characterization,kagan1973characterization}:
every non-Gaussian law has a nonzero rotational response.

\begin{lemma}[Gaussian characterization by the rotational Fourier defect]
\label{lem:fourier-ode}
For the product of centered complex exponentials defined above,
\begin{equation}
 R(\xi,\eta)
 =\left.\frac{d}{d\theta}
       \E W_{\xi,\eta}(O_\theta(X_1,X_2))\right|_{\theta=0}
   =\eta\phi'(\xi)\phi(\eta)
          -\xi\phi(\xi)\phi'(\eta).
 \label{eq:rotational-fourier-response}
\end{equation}
Moreover,
\begin{equation}
 R\equiv0
 \quad\Longleftrightarrow\quad
 \Law(X)\text{ is a centered Gaussian, possibly degenerate}.
 \label{eq:fourier-gaussian-characterization}
\end{equation}
Equivalently, all pairs of values of the weighted Fourier jet
\[
 u\longmapsto\bigl(u\phi(u),\phi'(u)\bigr)
\]
are linearly dependent if and only if \(\Law(X)\) is a centered Gaussian,
possibly degenerate.
\end{lemma}
\begin{proof}
Since $\E|X|<\infty$, we have $\phi\in C^1(\R)$ and
$\phi'(0)=i\E X=0$.  After rotation, the joint characteristic function is
\[
 \phi(\xi\cos\theta+\eta\sin\theta)
 \phi(-\xi\sin\theta+\eta\cos\theta).
\]
The two marginal terms in the centered product have zero derivative at
$\theta=0$, again by $\phi'(0)=0$; differentiating the displayed product
gives \eqref{eq:rotational-fourier-response}.

Now suppose $R\equiv0$.  Choose $\eta_0\ne0$ with
$\phi(\eta_0)\ne0$, possible because $\phi(0)=1$.  Then, for every $\xi$,
\[
 \phi'(\xi)=c\xi\phi(\xi),\qquad
 c=\frac{\phi'(\eta_0)}{\eta_0\phi(\eta_0)},\qquad
 \phi(0)=1.
\]
Uniqueness for this global linear ODE yields
$\phi(\xi)=e^{c\xi^2/2}$.  The identities
$\phi(-\xi)=\overline{\phi(\xi)}$ and $|\phi(\xi)|\le1$ imply respectively
$c\in\R$ and $c\le0$.  Thus $c=-\sigma^2$ for some $\sigma\ge0$, so $X$ is
centered Gaussian (the case $\sigma=0$ is the point mass at zero).
Conversely, direct substitution shows that every such Gaussian has
$R\equiv0$.
\end{proof}

\subsection{The determinant representation and finite detection}
\label{sec:fourier-determinant}

The determinant form supplies the finite frequency pair used to construct the
dual moments in \cref{sec:fourier-dual-moments}.

Consider the curve
\[
 u\longmapsto\bigl(u\phi(u),\phi'(u)\bigr)\in\mathbb C^2,
\]
described as the weighted Fourier jet.  Its values satisfy
\begin{equation}
 \det\!\begin{pmatrix}
 \xi\phi(\xi)&\phi'(\xi)\\
 \eta\phi(\eta)&\phi'(\eta)
 \end{pmatrix}
 =-R(\xi,\eta).
 \label{eq:fourier-wedge}
\end{equation}
The weighting by \(u\) is essential: for a nondegenerate Gaussian law, the
unweighted curve \(u\mapsto(\phi(u),\phi'(u))\) is not contained in a fixed
line.

For a non-Gaussian law, continuity supplies distinct frequencies and a compact
detecting neighborhood, independent of the ambient dimension; the diagonal
\(\xi=\eta\) never detects the law.

The same product of centered complex exponentials gives a local Hellinger
lower bound for one rotation.
Choose a phase $\alpha$ with
$\operatorname{Re}(e^{-i\alpha}R(\xi,\eta))>0$ and apply
\[
 |\E_\lambda F-\E_\nu F|
 \le h(\lambda,\nu)
             \{2\E_\lambda F^2+2\E_\nu F^2\}^{1/2}
\]
to $F=\operatorname{Re}(e^{-i\alpha}W_{\xi,\eta})$.
It follows that
$h((O_\theta)_\#\Law(X)^{\otimes2},
\Law(X)^{\otimes2})\ge c|\theta|$ for small
$|\theta|$.
A first moment suffices for this lower bound.
The dimension-free matrix estimates below require the additional
second-moment and Sobolev hypotheses.

\subsection{Dimension-free variance and affine normalization}
\label{sec:variance-and-affine-normalization}

The dimension-free variance bounds needed below are already contained in
\eqref{eq:centered-exponential-product-variance}; we record only how \(R\)
changes under affine normalization.

\paragraph{Centering, scaling, and phase.}
Let $Y$ have mean \(m\) and characteristic function \(\phi_Y\), and put
\(X=(Y-m)/\sigma\), where \(\sigma>0\).
For \(a=\xi/\sigma\) and \(b=\eta/\sigma\), the rotational Fourier defect
of $X$ satisfies
\begin{equation}
 \begin{aligned}
 R(\xi,\eta)
 &=e^{-i(m/\sigma)(\xi+\eta)}
 \bigl[b\phi_Y'(a)\phi_Y(b)-a\phi_Y(a)\phi_Y'(b)\\
 &\hspace{8em}+im(a-b)\phi_Y(a)\phi_Y(b)\bigr].
 \end{aligned}
 \label{eq:centering-phase}
\end{equation}
The bracketed term is the mean-corrected rotational response for two copies of
\(Y\).  Centering contributes the displayed phase, scaling changes the
frequencies, and the modulus at corresponding frequencies is unchanged.  For
a centered symmetric law the defect is real, but its phase at a single pair is
not a complete measure of skewness.

For the concrete comparison $Y\sim\operatorname{Gamma}(2,1)$,
$m=2$, $\sigma=\sqrt2$, and $(\xi,\eta)=(1,2)$, the expression in brackets in
\eqref{eq:centering-phase} equals $4/27$, so the law centered with variance one has
\[
 R(1,2)=\frac4{27}e^{-3i\sqrt2}.
\]
Thus the defect can remain explicit even for an asymmetric law with a support
boundary.

\subsection{The matrix Fisher operator}
\label{sec:fisher-operator}

The matrix Fisher operator separates diagonal scale, symmetric off-diagonal,
and skew-symmetric rotational directions.  Its rotational eigenvalue is
\(J-1\), which vanishes exactly in the Gaussian case.
The general interpretation of Fisher information as a metric on density
families is developed in~\cite{amari2000methods}; here the matrix decomposition
specifies the cost of each mixing direction.

Let $f$ be a probability density and write $g=\sqrt f$.
In this subsection we assume
\begin{equation}
 \begin{gathered}
 f>0\quad\text{Lebesgue-almost everywhere},\qquad
 g\in W^{1,2}(\R),\qquad xg'(x)\in L^2(\R),\\
 \E X=0,\qquad \E X^2=1,
 \qquad X\sim f(x)\,dx.
 \end{gathered}
 \label{eq:fisher-regularity}
\end{equation}
The derivatives in these assumptions are weak derivatives.
Define the location score $s=-f'/f=-2g'/g$ almost everywhere.  Then
\begin{equation}
 \E s(X)=0,\qquad \E[Xs(X)]=1,\qquad
 J:=\E s(X)^2=4\|g'\|_2^2.
 \label{eq:score-moments}
\end{equation}
Define the scale score and its second moment by
\begin{equation}
 d(x):=xs(x)-1,\qquad
 \kappa:=\E d(X)^2=4\|xg'\|_2^2-1.
 \label{eq:kappa}
\end{equation}
The two integration-by-parts identities in
\eqref{eq:score-moments} follow from these assumptions.
Indeed,
$f'=2gg'$ belongs to $L^1$, and $xf'$ belongs to $L^1$ by
Cauchy--Schwarz and the finite second moment.  Choose a smooth compactly
supported cutoff $\chi$ equal to one near zero.  Applying integration by
parts first to $\chi(x/L)$ and then to
$x\chi(x/L)$ gives the identities as $L\to\infty$.

For independent copies $X_1,\ldots,X_n$ of $X$, write
$\mathbf X^{(n)}=(X_1,\ldots,X_n)\sim\mu_n$.  For
$A\in M_n(\R)$, define the score of the pushforward by $I+tA$ at $t=0$ by
\begin{equation}
 S_A(\mathbf X^{(n)})=\sum_{i,j=1}^n A_{ij}
       \bigl(s(X_i)X_j-\delta_{ij}\bigr)
       =\sum_i A_{ii}d(X_i)+\sum_{i\ne j}A_{ij}s(X_i)X_j.
 \label{eq:matrix-score}
\end{equation}

\begin{theorem}[The matrix Fisher operator]
\label{thm:fisher-spectrum}
Under \eqref{eq:fisher-regularity}, let \(s\) and \(d\) be the location and
scale scores defined above, and define the matrix Fisher operator
$\mathcal I$ by
$\langle A,\mathcal I B\rangle_F=\E[S_AS_B]$.
Then
\begin{equation}
 \boxed{\mathcal I A=JA+A^*
                     +(\kappa-J-1)\diag A.}
 \label{eq:fisher-operator-spectrum}
\end{equation}
Here $\diag A$ is the matrix with diagonal entries $A_{ii}$ and zero
off-diagonal entries.  The diagonal, symmetric off-diagonal, and skew-symmetric matrix subspaces
are orthogonal for the Frobenius inner product and have eigenvalues
$\kappa$, $J+1$, and $J-1$, respectively.  Equivalently,
\begin{equation}
 \begin{aligned}
 \langle A,\mathcal I A\rangle_F
 &=\kappa\|\diag A\|_F^2
   +(J+1)\left\|\frac{A+A^*}2-\diag A\right\|_F^2\\
 &\quad +(J-1)\left\|\frac{A-A^*}2\right\|_F^2.
 \end{aligned}
 \label{eq:fisher-decomposition}
\end{equation}
The eigenvalue multiplicities are, respectively,
\begin{equation}
 n,\qquad n(n-1)/2,\qquad n(n-1)/2.
 \label{eq:fisher-multiplicities}
\end{equation}
Moreover, $\kappa>0$ and $J\ge1$.
Here $J=1$ holds precisely for the
standard Gaussian density.  Thus, for a non-Gaussian $f$,
$\min\{\kappa,J-1\}>0$ and
\begin{equation}
 \begin{gathered}
 \min\{\kappa,J-1\}\|A\|_F^2\le\langle A,\mathcal I A\rangle_F
 \le C_f\|A\|_F^2,\\
 C_f:=\max\{\kappa,J+1\}.
 \end{gathered}
 \label{eq:fisher-ellipticity}
\end{equation}
\end{theorem}

\begin{proof}
The variables $d(X_i)$ are centered and independent.  Their covariances with
the off-diagonal terms vanish, while two off-diagonal terms have nonzero
covariance only when their ordered index pairs coincide or are reversed; the
corresponding covariances are $J$ and $(\E Xs(X))^2=1$.  Hence
\begin{equation}
 \E S_A^2
 =\kappa\sum_i A_{ii}^2+
   \sum_{i<j}\{J(A_{ij}^2+A_{ji}^2)+2A_{ij}A_{ji}\}.
 \label{eq:fisher-pair-form}
\end{equation}
Polarization gives \eqref{eq:fisher-operator-spectrum}; resolving each
off-diagonal pair into its symmetric and skew-symmetric parts gives
\eqref{eq:fisher-decomposition} and the stated multiplicities.

Cauchy--Schwarz applied to $\E Xs(X)=1$ gives $J\ge1$.
Equality implies $s(X)=X$ almost surely, hence
$g'=-xg/2$ almost everywhere and $f$ is the standard Gaussian.
If $\kappa=0$, then $2xg'+g=0$ almost everywhere on each of the two
half-lines.  The absolutely continuous solutions there are constant
multiples of $|x|^{-1/2}$, which cannot give a positive probability
density.  This proves $\kappa>0$ and the asserted bounds.
\end{proof}

\subsection{Hellinger paths, polar decomposition, and symmetric perturbations}
\label{sec:hellinger-paths}

For orientation, write the left polar decomposition of an invertible matrix
as \(T=BO\), where \(B=(TT^*)^{1/2}\) is positive and \(O\) is
orthogonal.  The covariance defect is \(TT^*-I=B^2-I\), so a
finite second moment detects the symmetric factor.  The skew tangent of the
orthogonal factor has Fisher eigenvalue \(J-1\), which is positive precisely
off the Gaussian law.  The full matrix Fisher operator decomposition above packages
these two mechanisms without requiring a nonlinear separation of \(B\) and
\(O\) at every step.

The square-root formulation is the one used in quadratic-mean
differentiability~\cite[Section~7.2]{vandervaart1998asymptotic}.
Logarithmic derivatives and the relation between differentiation of measures
and quasi-invariance are treated in
\cite{bogachev1997differentiable,bell1985quasi}.
For the present linear transformations, the half-density action below gives
a path bound whose constant does not grow with the number of coordinates.

For $A\in\GL(n,\R)$ define a unitary operator on
$L^2(\R^n,dx)$ by
\[
 (U_Av)(y)=|\det A|^{-1/2}v(A^{-1}y),
 \qquad
 \Omega_n(x)=\prod_{i=1}^n g(x_i).
\]
Thus $U_A\Omega_n$ is the square root of the Lebesgue density of
$A_\#\mu_n$.  For invertible $A,B$,
\[
 h(A_\#\mu_n,B_\#\mu_n)=\|U_A\Omega_n-U_B\Omega_n\|_2.
\]
For a fixed real matrix $A$, the generator of $U_{\exp(tA)}$ is
\begin{equation}
 G_Av=-\tfrac12\operatorname{tr}(A)v-(Ax)\cdot\nabla v,
 \qquad G_A\Omega_n=\tfrac12S_A\Omega_n.
 \label{eq:half-density-generator}
\end{equation}
These identities hold in $L^2$.
To justify them under \eqref{eq:fisher-regularity}, observe that
$x_j\partial_i\Omega_n\in L^2$ for all $i,j$: the diagonal case
uses $xg'\in L^2$, and the off-diagonal case uses $g'\in L^2$
and $\E X^2=1$.  Cutoff and mollification approximate
$\Omega_n$ in the corresponding weighted Sobolev graph norm, where
the usual change-of-variables differentiation is valid.

\begin{lemma}[A dimension-free Hellinger path bound]
\label{lem:finite-hellinger-upper}
If $A_t$, $0\le t\le1$, is a continuously differentiable path in
$\GL(n,\R)$, then $t\mapsto U_{A_t}\Omega_n$ is continuously
differentiable in $L^2$ and
\begin{equation}
 \begin{aligned}
 \frac d{dt}U_{A_t}\Omega_n
 &=U_{A_t}G_{A_t^{-1}\dot A_t}\Omega_n,\\
 \left\|\frac d{dt}U_{A_t}\Omega_n\right\|_2
 &=\frac12\sqrt{\left\langle A_t^{-1}\dot A_t,
                  \mathcal I(A_t^{-1}\dot A_t)\right\rangle_F}.
 \end{aligned}
 \label{eq:finite-hellinger-path-derivative}
\end{equation}
Consequently,
\begin{equation}
 h((A_1)_\#\mu_n,(A_0)_\#\mu_n)
 \le\frac{\sqrt{C_f}}2
       \int_0^1\|A_t^{-1}\dot A_t\|_F\,dt.
 \label{eq:finite-hellinger-path}
\end{equation}
In particular, if $\|B\|_{\mathrm{op}}\le r<1$, then
\begin{equation}
 h((I+B)_\#\mu_n,\mu_n)
 \le\frac{\sqrt{C_f}}{2(1-r)}\|B\|_F.
 \label{eq:finite-hellinger-upper}
\end{equation}
All constants are independent of $n$.
\end{lemma}

\begin{proof}
The identity $U_AU_B=U_{AB}$ and \eqref{eq:half-density-generator} give the
derivative formula.  Unitarity and \eqref{eq:fisher-ellipticity} give its
norm; integration proves \eqref{eq:finite-hellinger-path}.  For the last
bound take $A_t=I+tB$ and use
$\|A_t^{-1}\|_{\mathrm{op}}\le(1-r)^{-1}$.
\end{proof}

\subsection{Finite real Fourier sums with dual moments}
\label{sec:fourier-dual-moments}

Assume in addition that \(f\) is non-Gaussian.  The score \(S_A\) is the correct
tangent, but it is not a usable dimension-free test after the law moves:
differentiating a score-based test would require regularity and integrability
not assumed here.  We instead need, for every matrix direction, a test built
from bounded one-coordinate functions, with exact duality against the score and
with \(L^2\)-norm bounded by a constant times the Frobenius norm of that
direction, uniformly in the dimension.  The rank-two
obstruction in \cref{lem:fourier-ode} produces the two off-diagonal coordinate
functions; a separate small-frequency argument using the second moment
produces the diagonal coordinate function.  All three are finite real Fourier
sums at frequencies depending on \(f\) but not on the matrix dimension.

\begin{lemma}[Finite real Fourier sums with dual moments]
\label{lem:fourier-dual-moments}
There exist centered finite real Fourier sums $p,q,r$ such that
\begin{equation}
 \begin{array}{c|cc}
       &s&x\\ \hline
 p     &1&0\\
 q     &0&1
 \end{array}
 \qquad\text{and}\qquad \E[r(X)d(X)]=1.
 \label{eq:fourier-dual-moments}
\end{equation}
The entries in the table denote expectations of the products.
For
\begin{equation}
 H_A(\mathbf X^{(n)}):=\sum_i A_{ii}r(X_i)
              +\sum_{i\ne j}A_{ij}p(X_i)q(X_j),
 \label{eq:fourier-functions}
\end{equation}
one has, for all real matrices $A,B$ of the same size,
\begin{equation}
 \E[H_AS_B]=\langle A,B\rangle_F,
 \qquad \|H_A\|_{L^2(\mu_n)}\le C_{\mathrm{dual}}\|A\|_F,
 \label{eq:matched-fourier-pairing}
\end{equation}
where $C_{\mathrm{dual}}$ is independent of $n$.
\end{lemma}

\begin{proof}
On the real space of centered finite real Fourier sums, consider
\[
 L_1(u)=\E u'(X)=\E[u(X)s(X)],
 \qquad L_2(u)=\E[Xu(X)].
\]
The integration-by-parts identity is valid because $u,u'$ are bounded.  If
$u\mapsto(L_1(u),L_2(u))$ had rank at most one, then after complexification
its determinant on the centered exponentials of frequencies $\xi,\eta$
would vanish.  Since
$L_1(e^{i\xi x})=i\xi\phi(\xi)$ and
$L_2(e^{i\xi x})=-i\phi'(\xi)$, this would give
\[
 \eta\phi'(\xi)\phi(\eta)
       -\xi\phi(\xi)\phi'(\eta)=0
 \qquad(\xi,\eta\in\R).
\]
By \cref{lem:fourier-ode}, this contradicts non-Gaussianity.  Thus the moment
map has rank two, and inverting the moment matrix of two Fourier sums gives
$p,q$ with the required dual moments.

For $\xi\ne0$ let
$r_0(x)=\cos(\xi x)-\E\cos(\xi X)$.  Dominated convergence
gives
\[
 \frac{\E[Xr_0'(X)]}{\xi^2}
 =-\E\left[X\frac{\sin(\xi X)}\xi\right]
 \longrightarrow-\E X^2=-1
 \qquad(\xi\to0).
\]
For small nonzero $\xi$ the denominator is nonzero, so
$r=r_0/\E[Xr_0'(X)]$ satisfies, by cutoff integration by parts,
$\E[r(X)d(X)]=\E[Xr'(X)]=1$.

Centering kills all diagonal--off-diagonal cross terms.  Coincident ordered
pairs contribute $\E[ps]\E[qX]=1$, whereas reversed pairs contribute
$\E[pX]\E[qs]=0$; hence $\E[H_AS_B]=\langle A,B\rangle_F$.
Finally, independence gives, for centered $a,b\in L^2(f)$,
\begin{equation}
 \left\|\sum_{i\ne j}A_{ij}a(X_i)b(X_j)\right\|_2
 \le\sqrt2\,\|a\|_2\|b\|_2\|A\|_F.
 \label{eq:quadratic-chaos-bound}
\end{equation}
Combining this with orthogonality of the centered diagonal terms proves the
uniform norm bound.
\end{proof}

The test functions \(H_A\) will be paired with the Hellinger tangent along a matrix
path.  To keep that pairing uniformly controlled after moving away from the
identity, we need one uniform derivative bound for \(H_A\Omega_n\).  A naive
term count grows with \(n\), while a direct H\"older estimate asks for
integrability that is unavailable.  The next lemma isolates this non-formal
step.  Its proof groups terms by their index-coincidence patterns, turning every
surviving coefficient array into a Frobenius-controlled contraction.  In
particular, it uses no fourth moment of \(X\) or higher moment of the score.

\begin{lemma}[Generator bound for the functions $H_A$]
\label{lem:fourier-generator-bound}
Let $p,q,r$ be any fixed centered bounded continuously differentiable
functions with bounded derivatives.  Define $H_A$ by
\eqref{eq:fourier-functions}.  There is a constant $C_{\mathrm{gen}}$,
depending on these functions and on $J,\kappa$, such that
\begin{equation}
 \|G_B(H_A\Omega_n)\|_{L^2(dx)}
 \le C_{\mathrm{gen}}\|A\|_F\|B\|_F
 \qquad(A,B\in M_n(\R),\ n\ge1).
 \label{eq:coordinate-function-generator-bound}
\end{equation}
\end{lemma}

\begin{proof}
For fixed $n$, both $H_A$ and its first derivatives are bounded.
Consequently
\[
 x_j\partial_i(H_A\Omega_n)
 =H_Ax_j\partial_i\Omega_n
       +x_j(\partial_iH_A)\Omega_n\in L^2(dx)
\]
by the weighted Sobolev regularity of $\Omega_n$ and
$\E X^2=1$.  The cutoff and mollification argument used for
\eqref{eq:half-density-generator} therefore puts $H_A\Omega_n$
in the domain of $G_B$.
The product rule gives
\begin{equation}
 G_B(H_A\Omega_n)
 =\left\{\tfrac12H_AS_B-(Bx)\cdot\nabla H_A\right\}\Omega_n.
 \label{eq:coordinate-function-generator-product}
\end{equation}
We bound the two terms uniformly in $n$ by grouping indices according to
their coincidences.

We will repeatedly use the following elementary consequence of
independence.  If $u_1,\ldots,u_m$ are centered $L^2(f)$ functions,
and a coefficient tensor $C$ is indexed by ordered tuples of
\emph{distinct} coordinates, then
\begin{equation}
 \left\|\sum_{i_1,\ldots,i_m}^{\ne}
 C_{i_1\ldots i_m}\prod_{a=1}^m u_a(X_{i_a})\right\|_2
 \le\sqrt{m!}\,\|C\|_{\ell^2}
                     \prod_{a=1}^m\|u_a\|_2.
 \label{eq:canonical-chaos-bound}
\end{equation}
Indeed, after squaring, a mixed term vanishes unless the two tuples have the
same index set; summing the at most $m!$ permutations and applying
Cauchy--Schwarz proves the estimate.  We use it only for $m\le4$.

First consider $H_AS_B$.  For a vector $a$ and a zero-diagonal
matrix $M$, abbreviate
\[
 D_a(u)=\sum_i a_i u(X_i),\qquad
 Q_M(u,v)=\sum_{i\ne j}M_{ij}u(X_i)v(X_j).
\]
The product $H_AS_B$ is the sum of the four products obtained by
multiplying $D_{\mathbf a}(r)+Q_{A_0}(p,q)$
and $D_{\mathbf b}(d)+Q_{B_0}(s,x)$, where
$\mathbf a=(A_{ii})_i$, $\mathbf b=(B_{ii})_i$,
$A_0=A-\diag A$, and
$B_0=B-\diag B$.
Every local product on a shared coordinate is a bounded function times one of
$d,s,x$, hence lies in $L^2(f)$; split it into its mean and centered part.
For $D_a(u)D_b(v)$, the distinct-index tensor, diagonal Hadamard vector, and
total contraction all have norm at most $\|a\|_2\|b\|_2$.  For
$D_a(u)Q_M(v,w)$, the three coincidence patterns give $a\otimes M$, a row
or column multiplication of $M$ by $a$, and $Ma$ or $M^*a$; all have norm at
most $\|a\|_2\|M\|_F$.  The reversed product is identical.

It remains only to record the contractions for
\[
 Q_M(u,v)Q_N(w,z)
 =\sum_{i\ne j,\,k\ne l}M_{ij}N_{kl}
                         u_i v_j w_k z_l,
\]
where $u,v$ are bounded and $w,z\in L^2$.  Four distinct indices give
$M\otimes N$.  For the representative one-coincidence pattern $i=k$, the
centered local product has coefficient norm
\[
 \left\{\sum_i\left(\sum_j M_{ij}^2\right)
                    \left(\sum_l N_{il}^2\right)\right\}^{1/2}
 \le\|M\|_F\|N\|_F.
\]
Its mean contraction is the off-diagonal part of $M^*N$; the other three
one-coincidence patterns are transposes of this one.  With both indices
shared, the coefficients are $M_{ij}N_{ij}$ or $M_{ij}N_{ji}$ and their row,
column, or total contractions.  For example,
\[
 \sum_i\left|\sum_j M_{ij}N_{ij}\right|^2
 \le\sum_i\left(\sum_j M_{ij}^2\right)
            \left(\sum_j N_{ij}^2\right)
 \le\|M\|_F^2\|N\|_F^2,
\]
and Cauchy--Schwarz gives the same bound for every transpose, Hadamard, and
total contraction.  These patterns exhaust all coincidences.  Applying
\eqref{eq:canonical-chaos-bound} to the centered pieces proves
\begin{equation}
 \|H_AS_B\|_2\le C\|A\|_F\|B\|_F.
 \label{eq:coordinate-function-score-bound}
\end{equation}

For $(Bx)\cdot\nabla H_A$, the diagonal part is
\[
 \sum_{i,k}A_{ii}B_{ik}r'(X_i)X_k.
\]
For $i\ne k$, centering $r'$ produces the matrix $(\diag A)B$ and the
vector $B^*\mathbf a-\mathbf a\mathbin{\odot}\mathbf b$; for $i=k$, center
$xr'(x)$.  Their coefficient norms are at most
$2\|A\|_F\|B\|_F$.

Let $M=A_0$.  The derivative of the
off-diagonal part is the sum of
\[
 F(M,B;p',q):=
 \sum_{i\ne j,k}M_{ij}B_{ik}p'(X_i)q(X_j)X_k
\]
and $F(M^*,B;q',p)$.  For $k\notin\{i,j\}$, centering $p'$ gives a
three-index coefficient of norm at most
\[
 \left\{\sum_i\left(\sum_jM_{ij}^2\right)
                    \left(\sum_kB_{ik}^2\right)\right\}^{1/2}
 \le\|M\|_F\|B\|_F.
\]
and the mean gives $M^*B-M^*\diag B$, of norm at most
$2\|M\|_F\|B\|_F$.  For $k=i$, centering $xp'$ gives the matrix
$(B_{ii}M_{ij})$ and vector $M^*\mathbf b$; for $k=j$, the contractions are
$M\mathbin{\odot}B$ and its row, column, and total sums.  The preceding
Hadamard estimates bound all of them by $\|M\|_F\|B\|_F$.  The local
functions lie in $L^2$ because $p',q$ are bounded and $\E X^2=1$.
The same argument applies to $F(M^*,B;q',p)$, so
\eqref{eq:canonical-chaos-bound} yields
\[
 \|(Bx)\cdot\nabla H_A\|_2\le C\|A\|_F\|B\|_F.
\]
Together with \eqref{eq:coordinate-function-score-bound} and
\eqref{eq:coordinate-function-generator-product}, this proves the lemma.
\end{proof}

We can now close the local argument.
The test functions \(H_A\) give a linear functional that sees the
first-order displacement, while the generator bound controls the error made
along the exponential path.
This turns the nonzero rotational Fourier defect into a uniform Hellinger lower bound
for every small matrix perturbation.

\begin{proposition}[Uniform local Hellinger coercivity]
\label{prop:local-fisher-hellinger}
For every non-Gaussian density satisfying
\eqref{eq:fisher-regularity}, there are constants
$c,C,\delta>0$, independent of $n$, such that
\begin{equation}
 c\|B\|_F\le h((I+B)_\#\mu_n,\mu_n)
             \le C\|B\|_F
 \qquad\text{whenever }\|B\|_F\le\delta.
 \label{eq:local-matrix-hellinger}
\end{equation}
The lower bound can be obtained using the fixed finite set of
coordinate frequencies from \cref{lem:fourier-dual-moments}.
\end{proposition}

\begin{proof}
Fix the finite real Fourier sums from \cref{lem:fourier-dual-moments}.
In exponential coordinates, unitarity and the group derivative give
\begin{align*}
 \langle H_A\Omega_n,(U_{\exp A}-I)\Omega_n\rangle
 &=\int_0^1
 \langle U_{\exp(-tA)}(H_A\Omega_n),G_A\Omega_n\rangle\,dt.
\end{align*}
The value obtained by replacing $U_{\exp(-tA)}(H_A\Omega_n)$ with
$H_A\Omega_n$ is $\frac12\E H_AS_A=\frac12\|A\|_F^2$.  Moreover,
$\|(U_{\exp(-tA)}-I)v\|_2\le t\|G_Av\|_2$ on the generator domain.
Therefore \cref{lem:fourier-generator-bound} and
\eqref{eq:fisher-ellipticity} yield
\[
 \left|\langle H_A\Omega_n,(U_{\exp A}-I)\Omega_n\rangle
                       -\tfrac12\|A\|_F^2\right|
 \le\frac{C_{\mathrm{gen}}\sqrt{C_f}}4\|A\|_F^3.
\]
For sufficiently small $\|A\|_F$, Cauchy--Schwarz and
\eqref{eq:matched-fourier-pairing} thus give
\[
 h((\exp A)_\#\mu_n,\mu_n)
 \ge\frac{1}{4C_{\mathrm{dual}}}\|A\|_F.
\]
For $\|B\|_{\mathrm{op}}<1/2$, the real power-series logarithm
$A=\log(I+B)$ exists and satisfies
\[
 \|A-B\|_F
 \le\sum_{k\ge2}\frac{\|B\|_{\mathrm{op}}^{k-1}}{k}\|B\|_F
 \le C\|B\|_{\mathrm{op}}\|B\|_F.
\]
Thus $\|A\|_F$ and $\|B\|_F$ are uniformly comparable in a fixed small
chart, proving the lower bound.  The upper bound is
\cref{lem:finite-hellinger-upper}.
\end{proof}

The half-density pairing supplies the uniform local certificate; tail symmetry
in \cref{sec:necessity} applies it to arbitrary output marginals without a
separate mixed-law variance bound for \(H_A\).

\subsection{Uniform coercivity modulo the exact linear stabilizer}
\label{sec:global-coercivity}

The following global form is proved in \cref{sec:globalization-proof}, after
the infinite-dimensional classification supplies the compactness input; it is
not used in that classification.

\begin{theorem}[Dimension-free global Hellinger coercivity]
\label{thm:global-coercivity}
Assume \eqref{eq:main-assumptions}.  Let \(\mathcal P_n\) denote the allowed
signed-permutation matrices in dimension \(n\), and set
\begin{equation}
 d_n(T)=\min_{Q\in\mathcal P_n}\|T-Q\|_F.
 \label{eq:distance-stabilizer}
\end{equation}
For every \(0<a\le b<\infty\)
there are constants \(c=c(f,a,b)>0\) and \(C=C(f,a,b)<\infty\) such
that, for every \(n\) and every \(T\in\GL(n,\R)\) whose singular
values lie in \([a,b]\),
\begin{equation}
 c\min\{d_n(T)^2,1\}
 \le h(\mu_n,T_\#\mu_n)^2
 \le C\min\{d_n(T)^2,1\}.
 \label{eq:global-coercivity}
\end{equation}
\end{theorem}
